\section{Discussion}
\label{sec:discussion}
\noindent\textbf{Multi-Agent Systems.}
\project{} currently targets tool invocations made through an LLM orchestration framework. However, emerging multi-agent systems may distribute tool use across multiple interacting agents. For example, one agent may invoke a tool and pass its result to another agent, which subsequently uses the result as an argument to a different tool. Supporting such systems requires extending \project{} from individual tool invocations to multi-agent execution paths. One possible approach is to construct a harness over the agent communication graph and fuzz both agent-to-agent messages and the tool invocations reachable from each agent. The tool-specific sanitizers generated by \project{} can remain unchanged, but must be applied at each tool boundary reached during the execution. Furthermore, testing multi-agent executions requires maintaining state across multiple agent and tool invocations. We leave support for multi-agent systems to future work.

\noindent\textbf{Coverage of Third-Party Components.}
Many tools are lightweight wrappers around external libraries, services, or applications. For example, a browser tool may contain only a small amount of orchestration code while most of its functionality executes inside Chromium. In such cases, coverage collected only from the tool implementation does not capture how extensively the underlying component is exercised. \project{} can be extended with component-specific coverage collectors that combine coverage from the tool implementation and its dependencies. For example, browser tools could additionally collect Chromium coverage, while tools invoking native libraries could collect coverage from instrumented library code. \project{} can then use the combined coverage when deciding whether a generated input exposes new behavior. Supporting arbitrary third-party components requires component-specific instrumentation interfaces, which we leave for future work.

\noindent\textbf{Stateful Tool Interactions.}
\project{} currently fuzzes tools by generating inputs for individual tool interfaces. However, some tool behavior depends on state established by previous invocations. For example, one tool invocation may create a resource that is subsequently modified, queried, or deleted by another invocation. Bugs involving such state may therefore require a sequence of otherwise valid tool calls. \project{} can be extended to generate multi-call harnesses by identifying relationships between tool arguments and outputs, such as a resource identifier returned by one tool and consumed by another. The fuzzer can then mutate both the arguments and ordering of calls while applying the existing semantic sanitizers at each tool boundary. Determining useful invocation sequences across large tool collections remains an open problem, and we leave stateful tool fuzzing to future work.
